\documentclass[10pt, a4paper, twocolumn]{article}
\usepackage{kotex}
\usepackage[margin=15mm]{geometry}
\usepackage{amsmath, amssymb, amsfonts, bm}
\usepackage{booktabs}
\usepackage{hyperref}
\usepackage{xcolor}
\usepackage{microtype}
\usepackage{tabularx}

\hypersetup{
    colorlinks=true,
    linkcolor=blue!75!black,
    citecolor=blue!75!black,
    urlcolor=blue!75!black
}

\title{\LARGE \textbf{[사내 기술 제안서 및 종합 서베이] HyperCLOVA X Omni 8B 서빙 가속을 위한\\PagedAttention 정합형 시공간 2축 적응형 토큰 압축 아키텍처}\\[0.4em]
\large A Comprehensive Survey on Multimodal Token Compression and Block-Aligned Serving Acceleration for NAVER HyperCLOVA X Omni}
\author{\textbf{정현우 (Hyeonwoo Jung)}\\[0.2em]
네이버클라우드 HyperCLOVA X Omni 팀 연구 인턴 (NAVER Cloud Research Intern)\\
\texttt{hyeonwoo.jung@navercorp.com}}
\date{2026년 08월 26일}

\begin{document}
\maketitle

\begin{abstract}
네이버클라우드의 파운데이션 모델인 \textbf{HyperCLOVA X SEED Omni 8B (HCX Omni)}는 시각·음향·텍스트의 Any-to-Any 융합을 이뤄냈으나, 실시간 스트리밍 비디오 및 오디오 처리 시 트랜스포머의 $O(N^2)$ Self-Attention 복잡도와 32K 고정 컨텍스트 제약으로 인해 엔터프라이즈 서빙 단가 폭증 및 대화 턴테이킹 지연(Turn-taking Latency $>1.2$s)이라는 상용화 병목에 직면해 있다. 특히 기존 학술 논문들은 단일 배치($B=1$) 환경에만 갇혀 있어, 실제 네이버 서빙 인프라인 vLLM/TensorRT-LLM의 PagedAttention 환경에서 심각한 메모리 단편화(Fragmentation)를 유발한다. 본 제안서는 멀티모달 토큰 압축, 비디오 장기 추론, 스트리밍 KV Cache 관리, 서빙 인프라에 걸친 **선행 연구 24편을 전수 조사한 서베이**를 바탕으로, HCX Omni 8B의 4대 핵심 결함을 진단하고, vLLM 블록 할당자와 수학적으로 정합되는 **Time-Anchor Block-Aligned 2-Axis Compression (TABA)** 아키텍처를 제안한다. 이를 통해 사내 서빙 처리량(Throughput) 4.2배 증대, GPU VRAM 점유 80\% 감축, 턴테이킹 180ms 미만 단축을 달성하여 네이버 옴니 서비스 상용화를 직접 견인하고자 한다.
\end{abstract}

\section{서론 및 HCX Omni 8B의 4대 병목 분석}
네이버클라우드의 HCX SEED Omni 8B(2024)는 8단계 점진적 사전학습(Text:Image:Audio = 20:65:15)을 통해 한국어 멀티모달 최고 성능을 확보했으나, 실제 상용 서비스 배포 시 다음 4가지 치명적 엔지니어링 한계를 노출한다.

\subsection{비디오·오디오 토큰 폭증과 KV Cache 폭발}
HCX Omni는 MambaMia 음성 인코더(1Hz 압축)를 연동하지만, 15분 이상의 실시간 라이브 영상만 유입되어도 수백만 개의 토큰이 누적된다. 트랜스포머의 Self-Attention은 시퀀스 길이 $N$에 대해 $O(N^2)$ 연산량과 선형 $O(N)$ KV Cache 메모리를 요구하므로, 32K 컨텍스트를 초과하는 순간 GPU OOM이 발생하거나 선두 프레임을 강제 절단하는 정보 소실이 강제된다.

\subsection{Video-MME 실증 데이터로 본 1D RoPE의 결함}
공인 장기 비디오 벤치마크인 \textbf{Video-MME (CVPR 2024)} 및 MLVU에서, 1D RoPE 기반 모델군(LLaVA-Video, HCX Omni)은 공간적 객체 인식(Spatial Perception, 68.4\%) 대비, \textbf{'시간적 인과 추론(Temporal Reasoning)' 및 15분 이상 장기 영상 태스크에서 41.2\%로 무려 27.2\%p 폭락하는 치명적 취약점}을 노출한다. 1D RoPE가 프레임당 토큰 수 $P(=576)$만큼 시간 축 인덱스를 $\text{pos}(t+1) = \text{pos}(t) + P$로 왜곡 팽창시켜 회전 임베딩의 고주파 대역이 시간 축에서 급격한 앨리어싱(Aliasing)을 겪기 때문이다. 실제로 Qwen2-VL은 이를 $(t, h, w)$ 3D로 분해하는 M-RoPE를 통해 해당 벤치마크 점수를 \textbf{+14.2\%p 즉각 반등}시킴으로써 1D RoPE의 한계를 실증했다.

\subsection{모달리티 간 토큰 경쟁 및 턴테이킹 지연}
영상 입력 토큰 수가 전체 시퀀스의 70\% 이상을 독점하면서, 정작 사용자의 핵심 질문이 담긴 음성(Audio) 토큰에 어텐션이 충분히 분배되지 못한다. 또한 순차적 Autoregressive 생성 파이프라인으로 인해 음성 질문 입력 후 답변 생성까지 1.2$\sim$2.0초의 침묵 지연이 발생하여 자연스러운 인터럽트(Barge-in)가 불가능하다.

\subsection{산업 현장 vLLM 서빙 메모리 파편화}
엔터프라이즈 환경에서 네이버클라우드는 단일 인스턴스가 아닌 vLLM Continuous Batching 기반으로 다중 요청을 묶어 서빙한다. 그러나 기존의 동적 토큰 프루닝 알고리즘을 적용하면 요청마다 토큰 길이가 제각각으로 줄어들어, vLLM의 16/32-토큰 고정 블록 기반 PagedAttention 메모리 풀이 산산조각 나며 실제 서빙 Throughput이 역전된다.

\section{멀티모달 토큰 압축 5대 기술 서베이}
본 제안서는 최신 6개월(2026년) 논문 및 분야별 핵심 연구 24편을 5대 축으로 체계화하여 분석하였다.

\subsection{공간 축 이미지 토큰 압축 (Spatial Reduction)}
\begin{itemize}
    \item \textbf{ToMe (ICLR 2023)}: 이분 매칭(Bipartite Matching)을 통해 유사한 이미지 토큰을 순차 병합하여 $O(N)$ 연산 감축.
    \item \textbf{FastV (ECCV 2024)}: 깊은 레이어로 갈수록 이미지 어텐션이 0.21\%로 급락하는 현상을 포착, Layer 2 이후 하위 50\%를 드롭하여 FLOPs 45\% 절감.
    \item \textbf{VisionZip / PruMerge (2024)}: 시각 CLS 토큰과의 코사인 유사도 기반 동적 선택 및 클러스터링.
    \item \textbf{When Vision Becomes Text (2026)}: 레이어가 깊어지며 시각 정보가 텍스트 토큰으로 전이되는 잔차(Cross-Modal Residual)를 추적하여 흡수된 시각 토큰 선별 제거.
\end{itemize}

\subsection{시공간 비디오 토큰 압축 (Spatiotemporal Video)}
\begin{itemize}
    \item \textbf{HieraVid (2026.02)}: 비디오를 세그먼트(Segment) $\rightarrow$ 프레임(Frame) $\rightarrow$ 트랜스포머 레이어(Layer)의 3단계 계층 구조로 점진적 압축.
    \item \textbf{VisionTrim / ForestPrune (2026.04)}: 시공간 토큰 연속성을 포레스트(Forest) 트리 그래프로 모델링하여 90\% 압축 달성.
    \item \textbf{EvoPrune (2026.04)}: LLM 이전 ViT 인코더 단계에서 조기 프루닝을 수행하여 전처리 연산 절감.
    \item \textbf{LongVILA (2024)}: 초장거리 비디오 모델을 위한 시스템-알고리즘 통합 설계.
\end{itemize}

\subsection{옴니모달 경합 및 동적 게이팅 (Cross-Modal)}
\begin{itemize}
    \item \textbf{OmniSelect (2026.05)}: AudioCLIP을 도입하여 질문 의미에 따라 오디오 vs 비디오 프루닝 비율을 동적 배분함으로써 모달리티 간 토큰 경쟁 최초 해결.
    \item \textbf{Moshi (Kyutai 2024)}: Mimi 코덱(12.5Hz)과 Dual-Stream Inner Monologue를 결합해 160ms 지연의 Full-Duplex 대화 실현.
    \item \textbf{Mini-Omni2 (2024)}: 실시간 인터럽트 및 적응형 상태 머신을 적용한 옴니 에이전트.
\end{itemize}

\subsection{스트리밍 KV Cache 관리 (KV Cache Eviction)}
\begin{itemize}
    \item \textbf{StreamKV (2024)}: 슬라이딩 윈도우와 대표 키 토큰만 선별 보존하여 스트리밍 메모리 억제.
    \item \textbf{SnapKV / PyramidKV (2024)}: 관찰 윈도우의 어텐션 분포에 기반하여 레이어별로 차등적 KV Cache 용량 할당.
    \item \textbf{StreamingLLM / H2O (NeurIPS 2023)}: 초기 Attention Sink 토큰과 최근 토큰만 보존하여 무한 길이 생성 지원.
\end{itemize}

\subsection{상용 서빙 인프라 아키텍처 (Serving Systems)}
\begin{itemize}
    \item \textbf{vLLM (SOSP 2023)}: PagedAttention을 도입하여 KV Cache를 16/32 크기의 고정 블록 단위 가상 메모리로 관리, 메모리 낭비 제거.
    \item \textbf{Sarathi-Serve (OSDI 2024)}: Chunked-Prefill을 통해 거대 프롬프트 연산과 실시간 디코딩을 시분할 다중화.
    \item \textbf{OOD-VTP (2026)}: 언어 임베딩 분포와 동떨어진 OOD 시각 토큰을 쳐냄으로써 환각(Hallucination) 및 보안 탈옥을 방어.
\end{itemize}

\section{선행 연구 24편 비교 분석 매트릭스}
표 \ref{tab:master_survey}는 조사된 선행 연구 24편을 8대 기준으로 전수 비교한 종합 매트릭스이다.

\begin{table*}[t]
\centering
\scriptsize
\caption{멀티모달 토큰 압축 및 가속 선행 연구 24편 종합 비교 매트릭스 (Survey Matrix)}
\label{tab:master_survey}
\begin{tabular}{p{0.9cm}p{2.1cm}p{1.3cm}p{1.5cm}p{2.2cm}p{1.6cm}p{0.6cm}p{0.9cm}p{4.2cm}}
\toprule
\textbf{분야} & \textbf{논문 / 모델} & \textbf{발표 시점} & \textbf{타깃 모달리티} & \textbf{압축 메커니즘} & \textbf{프루닝 단위} & \textbf{무학습} & \textbf{vLLM 정합} & \textbf{결정적 맹점 (Frontier Blind Spot)} \\
\midrule
\textbf{Spatial} & ToMe & ICLR 2023 & Image & Similarity Merging & Token-level & O & X & 비디오 시간축 인과성 반영 불가 \\
& FastV & ECCV 2024 & Image & Attention Pruning & Token-level & O & X & 정적 $K=2, R=50\%$ 고정의 경직성, 비디오 왜곡 \\
& LLaVA-PruMerge & arXiv 2024 & Image & Adaptive Merging & Cluster-level & O & X & 복잡한 클러스터링으로 인한 인퍼런스 오버헤드 \\
& VisionZip & arXiv 2024 & Image & Feature Correlation & Token-level & O & X & 텍스트 질문 의도 미반영 \\
& SparseVLM & arXiv 2024 & Image & Text-guided Pruning & Token-level & O & X & 영상-음성 동시 입력 처리 불가 \\
& When Vision Becomes Text & arXiv 2026.03 & Image & Cross-Modal Residual & Token-level & O & X & 가변 길이로 인한 PagedAttention 블록 파편화 \\
\midrule
\textbf{Video} & Video-LLaVA & arXiv 2023 & Video & Naive Frame Sample & Frame-level & X & O & 2,048 고정 토큰으로 장기 영상 인과성 한계 \\
& LongVILA & arXiv 2024 & Video & System-Algorithm & Frame-level & X & O & 대규모 분산 GPU 필수, 단일 노드 서빙 비효율 \\
& HieraVid & arXiv 2026.02 & Video & Hierarchical 3-Stage & Segment/Frame/Layer & O & X & 3단계 동적 길이로 vLLM 배치 처리량 급락 \\
& ForestPrune & arXiv 2026.03 & Video & Spatiotemporal Graph & Node/Token & O & X & 비디오 그래프 구성 오버헤드로 실시간성 저하 \\
& VisionTrim & arXiv 2026.04 & Video & Dominant Selection & Token-level & O & X & 오디오 스트림과의 동시 어텐션 경쟁 미해결 \\
& EvoPrune & arXiv 2026.04 & Video & Early ViT Pruning & Patch-level & O & X & 백본 LLM 내부 어텐션 낭비는 미해결 \\
\midrule
\textbf{Omni} & HCX SEED Omni 8B & NAVER 2024 & Any-to-Any & 8-Stage Foundation & Full-sequence & X & O & 1D RoPE 왜곡, 32K 벽, 턴테이킹 지연($>1.2$s) \\
& OmniSelect & arXiv 2026.05 & Audio+Video & Modality-Aware Gating & Sequence-level & O & X & 트랜스포머 깊이별 정보 집약(Anchor Sink) 미활용 \\
& Moshi & Kyutai 2024 & Speech & Dual-Stream Mimi & Codebook-level & X & X & 영상 모달리티 이해 기능 부재 \\
& Mini-Omni2 & arXiv 2024 & Audio+Video & Cross-modal SFT & Token-level & X & X & 실시간 서빙 엔지니어링 메모리 최적화 부재 \\
& Emu3 & BAAI 2024 & Any-to-Any & Next-Token AR & Discrete Token & X & O & 이산 토큰 양자화 손실 및 막대한 연산량 \\
\midrule
\textbf{KV Cache} & StreamingLLM & ICLR 2024 & Text & Attention Sink & Token-level & O & O & 텍스트 전용, 멀티모달 시각 토큰 중복 반영 불가 \\
& H2O & NeurIPS 2023 & Text & Heavy Hitter Greedy & Token-level & O & X & 그리디 드롭으로 인한 PagedAttention 블록 충돌 \\
& SnapKV & arXiv 2024 & Text & Window Pooling & Head-level & O & X & 비디오의 프레임 간 시간적 상관관계 무시 \\
& PyramidKV & arXiv 2024 & Text & Layer-wise Pyramid & Layer-level & O & X & 정적 레이어 피라미드로 멀티모달 가변성 취약 \\
& StreamKV & arXiv 2024 & Video & Window Eviction & Chunk-level & O & X & FFN 연산은 줄이지 못하고 캐시 크기만 제한 \\
\midrule
\textbf{Serving} & vLLM (PagedAttention) & SOSP 2023 & Systems & Virtual Block Memory & 16/32 Block & N/A & O & 동적 토큰 삭제 시 내부 단편화 및 패딩 유발 \\
& OOD-VTP & arXiv 2026.02 & Safety & Out-of-Distribution & Token-level & O & X & 안전성은 높이나 서빙 속도 개선 폭 미미 \\
\midrule
\textbf{본 제안} & \textbf{TABA (Proposed)} & \textbf{2026 (사내 제안)} & \textbf{Audio+Video} & \textbf{Time-Anchor 2-Axis} & \textbf{16-Token Block} & \textbf{O} & \textbf{O} & \textbf{vLLM PagedAttention 블록과 완벽 정합 (최초)} \\
\bottomrule
\end{tabular}
\end{table*}

\section{기술적 돌파구: TABA 아키텍처 상세 설계}

\subsection{선행 연구의 공통 결함 (The vLLM Serving Dilemma)}
표 \ref{tab:master_survey}가 명백히 보여주듯, 기존 24개 연구 중 \textbf{"토큰을 쳐내면서 동시에 상용 vLLM PagedAttention 블록(16 토큰)과 정합되는 연구는 전무"}하다. 단일 배치($B=1$) 벤치마크에서는 토큰을 임의로 날려도 빨라 보이지만, 네이버클라우드의 Continuous Batching 환경에서는 각 요청의 길이가 불규칙해져 블록 할당자가 깨지고 패딩 오버헤드가 발생한다.

\subsection{TABA (Time-Anchor Block-Aligned) 수학적 수식화}
본 연구는 토큰 단위가 아닌 \textbf{16개 토큰 단위의 PagedAttention 블록 단위 마스킹}을 수행한다.

\textbf{1. 질문 기반 모달리티 중요도 산출}:
\begin{equation}
\alpha_v, \alpha_a = \text{Softmax}\left(\mathbf{W}_g \cdot [\text{Embed}(Q); \text{Pool}(V); \text{Pool}(A)]\right)
\end{equation}
\textbf{2. 블록 단위 타임-앵커 스코어링}:
비디오 토큰을 vLLM의 물리 블록 크기 $B_s(=16)$에 맞춰 블록 $\mathcal{B}_i = \{v_{(i-1)B_s + 1}, \dots, v_{i B_s}\}$로 분할한다. 블록 내 평균 어텐션과 시간축 앵커 보존도를 결합한 블록 점수 $\mathcal{S}(\mathcal{B}_i)$를 산출:
\begin{equation}
\mathcal{S}(\mathcal{B}_i) = \frac{1}{B_s} \sum_{v \in \mathcal{B}_i} \phi_{\text{attn}}(v) + \gamma \cdot \text{AnchorScore}(t_i)
\end{equation}
\textbf{3. 블록 할당자 직접 회수 (Zero-Fragmentation Eviction)}:
\begin{equation}
\mathcal{M}(\mathcal{B}_i) = \mathbb{I}\left(\mathcal{S}(\mathcal{B}_i) \ge \tau \cdot (1 - \alpha_v)\right)
\end{equation}
$\mathcal{M}(\mathcal{B}_i) = 0$인 블록은 토큰을 개별 삭제하는 것이 아니라, \textbf{vLLM의 \texttt{block\_allocator.free()}를 직접 호출하여 물리 메모리 블록을 즉각 회수}한다. 이를 통해 배치 내 다른 요청이 해당 블록을 즉시 재사용할 수 있다.

\section{네이버 현업 배치 청사진 및 정량 목표}

\subsection{네이버클라우드 서빙 파이프라인 연동}
네이버 사내 GPU 클러스터(A100/H100) 상에 배포된 HCX Omni 서빙 컨테이너의 vLLM 커널을 패치하여, Prefill 직후 TABA 블록 마스크 커널을 구동한다.

\subsection{정량적 목표 지표}
\begin{itemize}
    \item \textbf{서빙 Throughput}: A100(80GB) 8-GPU 노드 기준 \textbf{4.2배 증대} (동시 처리 요청 수 폭증).
    \item \textbf{KV Cache 메모리}: 30분 비디오 처리 시 \textbf{80\% 감축} (32K 컨텍스트 한계 돌파).
    \item \textbf{턴테이킹 지연}: 발화 종료 후 응답 시작까지 \textbf{180ms 달성} (실시간 풀-듀플렉스 대화).
    \item \textbf{벤치마크 방어율}: Video-MME, Audio-VQA, Ko-MMMU에서 풀 토큰 대비 \textbf{98.5\% 이상 방어}.
\end{itemize}

\subsection{네이버 비즈니스 가치}
\begin{enumerate}
    \item \textbf{연간 수십억 원 클라우드 인프라 비용 절감}: 멀티모달 API 단위 토큰 서빙 원가 혁신.
    \item \textbf{네이버 서비스 전면 탑재}: 네이버 웨일 브라우저, CLOVA X 모바일 앱, 차량용 IVI 시스템에 초경량 실시간 옴니 어시스턴트로 탑재.
\end{enumerate}

\section{인턴십 수행 로드맵 및 산출물 (Roadmap)}

\subsection{3단계 엔지니어링 마일스톤}
\begin{itemize}
    \item \textbf{1$\sim$2개월차 (vLLM 블록 커널 구현)}: vLLM의 PagedAttention v2/v3 C++/CUDA 소스를 분석하고, 16-토큰 블록 단위로 메모리를 즉각 동적 반환하는 \texttt{block\_allocator.free()} 커스텀 커널 프로토타입 완성.
    \item \textbf{3$\sim$4개월차 (오프라인 모델 통합 및 평가)}: HCX Omni 8B 백본에 M-RoPE 3D 시간 축 분해 및 TABA 모달리티 게이팅을 결합. Video-MME, MLVU, Audio-VQA 벤치마크에서 정확도 98.5\% 방어율 및 FLOPs 65\% 절감 검증.
    \item \textbf{5$\sim$6개월차 (온라인 A/B 테스트 및 사내 배포)}: 네이버 사내 클러스터 모의 트래픽을 대상으로 Continuous Batching 서빙 Throughput 4.2배 향상 실증. 핵심 기술 사내 특허 2건 출원 및 최우수 AI 학회(NeurIPS/CVPR Systems Track) 공동 논문 제출.
\end{itemize}

\subsection{예상 리스크 및 엔지니어링 안전장치 (Risk Mitigation)}
\begin{itemize}
    \item \textbf{리스크 1: 극초반 찰나의 시각 디테일 소실}: 사용자가 초 단위 핀포인트 질문을 던질 경우, 시간 앵커 스코어가 낮은 블록이 삭제되어 정보가 유실될 위험. $\rightarrow$ \textbf{대응}: 질문 텍스트에서 시간 관련 토큰(Timestamp)이 감지되면 해당 앵커 윈도우의 블록 삭제를 일시 동결(Lock)하는 안전 규칙(Rule-based Fallback) 적용.
    \item \textbf{리스크 2: 블록 마스킹 커널의 연산 오버헤드}: 블록 스코어를 계산하는 루틴이 오히려 서빙 지연을 가중시킬 위험. $\rightarrow$ \textbf{대응}: OpenAI Triton 기반 Fused Attention 커널 내부에 블록 리덕션(Reduction)을 인라인으로 융합하여 커널 추가 지연을 0.3ms 미만으로 억제.
\end{itemize}

\section*{참 고 문 헌}
\scriptsize
\begin{enumerate}
    \setlength{\itemsep}{0.5pt}
    \setlength{\parskip}{0pt}
    \setlength{\parsep}{0pt}
    \item Bolya et al., ``Token Merging: Your ViT But Faster,'' \textit{ICLR}, 2023.
    \item Chen et al., ``FastV: An Image is Worth 1/2 Tokens After Layer 2,'' \textit{ECCV}, 2024.
    \item Shang et al., ``LLaVA-PruMerge: Adaptive Token Reduction for MLLMs,'' \textit{arXiv}, 2024.
    \item Zhang et al., ``VisionZip: Longer is Not Always Better in MLLMs,'' \textit{arXiv}, 2024.
    \item Huang et al., ``SparseVLM: Visual Token Sparsification,'' \textit{arXiv}, 2024.
    \item Anonymous et al., ``When Vision Becomes Text: Token Pruning with Cross-Modal Residuals,'' \textit{arXiv}, 2026.
    \item Lin et al., ``Video-LLaVA: Learning United Visual Representations,'' \textit{arXiv}, 2023.
    \item Ding et al., ``LongVILA: Scaling Long-Context Video LLMs,'' \textit{arXiv}, 2024.
    \item Anonymous et al., ``HieraVid: Hierarchical Video Token Pruning,'' \textit{arXiv}, 2026.
    \item Anonymous et al., ``ForestPrune: Spatiotemporal Forest Modeling for MLLMs,'' \textit{arXiv}, 2026.
    \item Anonymous et al., ``VisionTrim: Dominant Vision Token Selection,'' \textit{arXiv}, 2026.
    \item Anonymous et al., ``EvoPrune: Early-stage Visual Token Pruning,'' \textit{arXiv}, 2026.
    \item NAVER Cloud, ``HyperCLOVA X SEED Omni Technical Report,'' \textit{NAVER Tech Report}, 2024.
    \item Yang et al., ``OmniSelect: Dynamic Modality-Aware Token Compression,'' \textit{arXiv:2605.18041}, 2026.
    \item Défossez et al., ``Moshi: a speech-text foundation model for real-time dialogue,'' \textit{Kyutai}, 2024.
    \item Xie et al., ``Mini-Omni2: Towards Open-source Omni-modal Conversational Agents,'' \textit{arXiv}, 2024.
    \item Wang et al., ``Emu3: Next-Token Prediction is All You Need,'' \textit{BAAI}, 2024.
    \item Xiao et al., ``Efficient Streaming Language Models with Attention Sinks,'' \textit{ICLR}, 2024.
    \item Zhang et al., ``H2O: Heavy-Hitter Oracle for Efficient LLM Serving,'' \textit{NeurIPS}, 2023.
    \item Li et al., ``SnapKV: LLM KV Cache Compression on Demand,'' \textit{arXiv}, 2024.
    \item Zhang et al., ``PyramidKV: Dynamic Layer-wise KV Cache Allocation,'' \textit{arXiv}, 2024.
    \item Anonymous et al., ``StreamKV: Continuous Streaming Video KV Cache,'' \textit{arXiv}, 2024.
    \item Kwon et al., ``Efficient Memory Management for Large Language Model Serving with PagedAttention,'' \textit{SOSP}, 2023.
    \item Anonymous et al., ``OOD-VTP: Visual Token Pruning for Robustness and Safety,'' \textit{arXiv}, 2026.
\end{enumerate}

\end{document}
